\title{Group Sequence Policy Optimization}

\begin{abstract}
We present Group Sequence Policy Optimization (GSPO), a reinforcement learning algorithm designed to provide stability, efficiency, and high performance in training large language models. Departing from prior methods that utilize token-level importance ratios, GSPO establishes importance ratios at the level of entire sequences and conducts sequence-wise clipping, reward calculation, and optimization. Our experiments show that GSPO delivers enhanced training speed and effectiveness relative to the GRPO algorithm, significantly stabilizes reinforcement learning for Mixture-of-Experts (MoE) architectures, and has the capacity to streamline RL system implementation. The advances achieved by GSPO have played a key role in driving notable enhancements in the latest Qwen3 models.
\end{abstract}

\section{Introduction}
\label{sec:intro}

Reinforcement learning (RL) has become a central approach in advancing the scalability of language models \citep{o1,dpsk-r1,qwq32b,qwen3}. Employing RL at a large scale enables language models to engage in more complex tasks—such as solving high-level mathematics and programming challenges—by facilitating more extensive and intricate reasoning abilities.

Achieving effective scaling of RL through increased computational resources critically depends on preserving stable and reliable training behavior. Nevertheless, leading RL algorithms such as GRPO \citep{grpo} suffer from pronounced stability challenges during the training of large-scale language models, frequently leading to drastic and permanent failures—referred to as model collapse \citep{qwen3, minimax-m1}. Such instability poses a significant barrier to advancing the capabilities of language models by means of further RL-based optimization.

This study reveals that the core instability of GRPO originates from the incorrect use and breakdown of importance sampling weights within its algorithm. Such misapplication leads to significant training noise with elevated variance, which intensifies as the response length grows and is exacerbated by the clipping strategy, eventually resulting in the collapse of the model.

In response to these fundamental challenges, we introduce \textbf{Group Sequence Policy Optimization (GSPO)}, an RL-based framework for optimizing large language models. GSPO’s principal contribution is a rigorously defined importance ratio rooted in sequence likelihood \citep{click}, thereby adhering closely to the foundational ideas of importance sampling. Moreover, GSPO derives normalized rewards by comparing the advantages of several responses generated for a single query, which promotes consistency between sequence-level reward assignment and the corresponding optimization objective.

Our experimental results reveal that GSPO markedly outperforms GRPO in terms of training stability, efficiency, and overall effectiveness. Notably, GSPO intrinsically addresses the instability issues commonly encountered during RL training with large Mixture-of-Experts (MoE) architectures, thus obviating the requirement for intricate stabilization techniques and potentially streamlining RL system design. The advantages offered by GSPO have played a crucial role in driving the remarkable gains observed in the most recent Qwen3 models. We anticipate that GSPO will serve as a reliable and extensible foundational approach, supporting the ongoing progress of large-scale RL training in language model development.

\section{Preliminaries}

\paragraph{Notation}
Throughout this work, we consider an autoregressive language model with parameters $\theta$, which serves as a policy, denoted by $\pi_\theta$. The symbol $x$ represents a query, while the collection of all queries is designated as $\mathcal{D}$. For any query $x$ and its associated response $y$, the probability of generating $y$ according to the policy $\pi_\theta$ is given by $\pi_\theta (y | x)=\prod_{t=1}^{|y|} \pi_\theta (y_t | x, y_{<t} )$, where $|y|$ indicates the total number of tokens in response $y$. Each pair $(x, y)$ consisting of a query and a response is evaluated by a verifier $r$ that assigns a reward $r(x, y)$ in the interval $[ 0, 1 ]$.

\paragraph{Proximal Policy Optimization (PPO)} PPO \citep{ppo} leverages data sampled from the previous policy $\pi_{\theta_\text{old}}$, enforcing policy updates to remain close to the old policy by means of a clipping strategy. More precisely, PPO formulates policy optimization around the following surrogate objective (for conciseness, we disregard the KL regularization term in what follows, as it lies outside the primary scope of this work):

\begin{align}
\mathcal{J}_\text{PPO}(\theta) = \mathbb{E}_{ x \sim \mathcal{D},\, y \sim \pi_{\theta_\text{old}}( \cdot | x) }
\left[ \frac{1}{|y|} \sum_{t=1}^{|y|} 
\min \left( w_{t}(\theta) \widehat{A}_{t},  \, \mathrm{clip} \left( w_{t}(\theta), 1 - {\varepsilon}, 1 + {\varepsilon}\right) \widehat{A}_{t} \right)
\right],
\end{align}

Here, the token $y_t$'s importance ratio is given by $
w_{t}(\theta) = \frac{ \pi_{\theta} (y_{t} | x, y_{<t}) }{ \pi_{\theta_\text{old}} (y_{t} | x,y_{<t})} $, where $\widehat{A}_{t}$ denotes the estimated advantage for $y_t$, computed via a separate value function. The parameter $\varepsilon$ specifies the clipping threshold applied to the importance ratios.

A primary difficulty with PPO in real-world applications stems from its significant dependence on the value model. Notably, the value model often matches the policy model in scale, resulting in substantial demands on both memory and computation. In addition, the algorithm's performance is closely tied to the precision of its value estimation. Although developing an accurate value model is fundamentally demanding, the challenges intensify when adapting the model to handle longer outputs and increasingly sophisticated tasks.

\paragraph{Group Relative Policy Optimization (GRPO)} GRPO \citep{grpo} eliminates reliance on a value function by evaluating the advantage of individual responses in the context of a group of responses corresponding to the same prompt. More precisely, GRPO seeks to optimize the following objective:

\begin{align}
\label{equ:grpo}
\mathcal{J}_\text{GRPO}(\theta) = \mathbb{E}_{ x \sim \mathcal{D},\, \{y_i\}_{i=1}^G \sim \pi_{\theta_\text{old}}( \cdot | x) }
\Bigg[ \frac{1}{G} \sum_{i=1}^{G} \frac{1}{|y_i|} \sum_{t=1}^{|y_i|} 
\min \Big( w_{i,t}(\theta) \widehat{A}_{i,t}, \,
\mathrm{clip} \big( w_{i,t}(\theta), 1 - {\varepsilon}, 1 + {\varepsilon}\big) \widehat{A}_{i,t} \Big)
\Bigg],
\end{align}

where $G$ is the number of generated responses to each query $x$ (i.e., the group size), and the importance ratio $w_{i,t}(\theta)$ and advantage $\widehat{A}_{i,t}$ of token $y_{i,t}$ are:

\begin{align}
    w_{i,t}(\theta) = \frac{ \pi_{\theta} (y_{i,t} | x, y_{i,<t}) }{ \pi_{\theta_\text{old}} (y_{i,t} | x, y_{i,<t}) }, \quad\ 
    \widehat{A}_{i,t} = \widehat{A}_i = \frac{ r(x, y_i) - \mathrm{mean}\left( \{ r(x, y_i) \}_{i=1}^G \right) }{ \mathrm{std}\left( \{ r(x, y_i) \}_{i=1}^G \right) },
\end{align}

respectively, where all the tokens in $y_i$ share the same advantage as $\widehat{A}_{i}$.

\section{Motivation}
\label{sec:motivation}

The increasing scale of models, the use of sparsity (such as in Mixture-of-Experts architectures), and longer response outputs require employing large rollout batch sizes to fully exploit hardware resources during reinforcement learning. In order to boost sample efficiency, it is common to divide a substantial rollout batch into several smaller mini-batches that are then used for gradient computations. However, this approach inherently creates an off-policy learning scenario, as the generated responses $y$ come from a previous policy $\pi_{\theta_\text{old}}$, rather than the policy $\pi_\theta$ that is currently being improved. As a result, techniques like PPO and GRPO employ a clipping strategy, which serves to exclude samples that are too far from the present policy, thereby ensuring the reliability of the gradient estimates.

Although techniques such as clipping are employed to address off-policy discrepancies, we observe a deeper problem within GRPO: \textit{the objective itself is ill-posed}. This issue is especially pronounced when applying GRPO to the training of large-scale models engaged in tasks that require generating extended responses, where it can result in severe model degeneration. The root cause of this ill-posedness lies in an incorrect use of importance sampling weights. In importance sampling, the underlying concept is to approximate the expectation of a function $f$ with respect to a target distribution $\pi_\text{tar}$ by appropriately reweighting data sampled from a behavior distribution $\pi_\text{beh}$:

\begin{align}
\mathbb{E}_{z \sim \pi_\text{tar}} \left[ f(z) \right] = \mathbb{E}_{z \sim \pi_\text{beh}} \left[ \frac{ \pi_\text{tar} (z) }{ \pi_\text{beh} (z)} \, f(z) \right].
\end{align}

A key aspect of this approach is that averaging across a large number of samples ($N \gg 1$) drawn from the behavior distribution $\pi_\text{beh}$ is essential. This ensures that the importance weight $\frac{ \pi_\text{tar} (z) }{ \pi_\text{beh} (z)}$ can adequately compensate for differences between the distributions.

On the other hand, GRPO incorporates the importance weight $\frac{ \pi_{\theta} (y_{i,t} | x, y_{i,<t}) }{ \pi_{\theta_\text{old}} (y_{i,t} | x,y_{i,<t})}$ at each individual token position $t$. This weighting is calculated using only a single observation $y_{i,t}$, drawn from the corresponding next-token distribution $\pi_{\theta_\text{old}}( \cdot | x, y_{i, <t})$. As a consequence, GRPO does not achieve effective correction for the distribution mismatch as intended. Rather, it generates substantial variance in the gradient estimates during training. This elevated noise can accumulate, especially with longer sequences, and its effect is further magnified when the clipping mechanism is used. Through empirical analysis, we have found that this can provoke catastrophic model collapse—a phenomenon from which recovery is rarely possible. In such cases, even restoring the model from earlier checkpoints, adjusting hyperparameters (such as modifying the clipping range), increasing sequence generation lengths, or altering the nature of RL queries fails to revive effective training.

The preceding analysis reveals a central flaw in the architecture of GRPO. Specifically, the ineffectiveness of token-level importance weighting highlights a crucial insight: \textbf{the optimization objective must be aligned with the granularity of the reward}. Given that rewards are distributed based on complete sequences, employing off-policy corrections for individual tokens proves inadequate. Consequently, we are prompted to abandon token-level objectives in favor of implementing importance weighting and conducting optimization processes at the \textit{sequence level} instead.

\section{Algorithm}

\subsection{GSPO: Group Sequence Policy Optimization}

Although the token-level importance weight $\frac{ \pi_{\theta} (y_{i,t} | x, y_{i,<t}) }{ \pi_{\theta_\text{old}} (y_{i,t} | x,y_{i,<t})}$ presents challenges within GRPO, our analysis reveals that, in language generation scenarios, the \textit{sequence-level} importance weight $\frac{ \pi_{\theta} (y | x) }{ \pi_{\theta_\text{old}} (y | x)}$ possesses a well-defined theoretical interpretation. Specifically, it quantifies the extent to which the response $y$—sampled from $\pi_{\theta_\text{old}} (\cdot | x)$—diverges from $\pi_{\theta} (\cdot | x)$. This directly connects with the sequence-level reward and offers a substantive measure for understanding and applying the clipping strategy.

Based on this straightforward observation, we propose the \textbf{Group Sequence Policy Optimization (GSPO)} algorithm. GSPO employs the following sequence-level optimization objective:

\begin{align}
\mathcal{J}_\text{GSPO} (\theta) =
\mathbb{E}_{ x \sim \mathcal{D},\, \{y_i\}_{i=1}^G \sim \pi_{\theta_\text{old}}( \cdot | x) }
\left[ 
\frac{1}{G} \sum_{i=1}^{G}
\min \Big( s_{i}(\theta)  \widehat{A}_{i}, \; \mathrm{clip} \big( s_{i}(\theta), 1 - {\varepsilon}, 1 + {\varepsilon} \big) \widehat{A}_{i} \Big) 
\right],
\label{equ:gspo}
\end{align}

where we adopt the group-based advantage estimation:

\begin{align}
\widehat{A}_{i} = \frac{r(x, y_i) - \mathrm{mean} \left( \{ r(x, y_i) \}_{i=1}^G \right) }{ \mathrm{std} \left( \{ r(x, y_i) \}_{i=1}^G \right) },
\end{align}

and define the importance ratio $s_{i}(\theta)$ based on sequence likelihood \citep{click}:

\begin{align}
s_{i}(\theta) = \left( \frac{ \pi_{\theta} (y_i | x) }{ \pi_{\theta_\text{old}} (y_i | x)} \right)^{\frac{1}{|y_i|}}
=
\exp \left( \frac{1}{|y_i|} \sum_{t=1}^{|y_i|} \log \frac{ \pi_{\theta} (y_{i,t} | x, y_{i,<t}) }{ \pi_{\theta_\text{old}} (y_{i,t} | x, y_{i,<t}) } \right).
\end{align}

Accordingly, GSPO implements clipping at the response level, rather than on individual tokens, to eliminate highly “off-policy” samples from gradient computation; this approach aligns with both sequence-level reward allocation and optimization procedures. Importantly, we employ length normalization in $s_{i}(\theta)$ to limit variance and to ensure $s_{i}(\theta)$ remains within a standardized numerical scale. Without this normalization, substantial changes in the probability of only a few tokens could cause significant volatility in the overall sequence-level importance ratio, necessitating different clipping ranges for responses of varying lengths. Additionally, it should be pointed out that the scale of the clipping ranges in GSPO is generally different—by orders of magnitude—from those used in prior algorithms such as GRPO, reflecting differences in how importance ratios are formulated.

\subsection{Gradient Analysis}
\label{subsec:gradient}

We can derive the gradient of the GSPO objective as follows (clipping is omitted for brevity):

\begin{align}
\nabla_{\theta} \mathcal{J}_\text{GSPO} (\theta)
=&\ 
\nabla_{\theta} \mathbb{E}_{ x \sim \mathcal{D},\, \{y_i\}_{i=1}^G \sim \pi_{\theta_\text{old}}( \cdot | x) }
\Bigg[ 
\frac{1}{G} \sum_{i=1}^{G} s_{i}(\theta) \widehat{A}_{i}
\Bigg] \\
= &\ 
\mathbb{E}_{ x \sim \mathcal{D},\, \{y_i\}_{i=1}^G \sim \pi_{\theta_\text{old}}( \cdot | x) }
\Bigg[ 
\frac{1}{G} \sum_{i=1}^{G}
s_{i}(\theta) \widehat{A}_{i}
\, \nabla_{\theta} \log s_{i}(\theta)
\Bigg] \\
=&\ 
\mathbb{E}_{ x \sim \mathcal{D},\, \{y_i\}_{i=1}^G \sim \pi_{\theta_\text{old}}( \cdot | x) }
\Bigg[ 
\frac{1}{G} \sum_{i=1}^{G} 
\left( \frac{ \pi_{\theta} (y_i | x) }{ \pi_{\theta_\text{old}} (y_i | x)} \right)^{\frac{1}{|y_i|}} \widehat{A}_{i} 
\, \frac{1}{|y_i|} \sum_{t=1}^{|y_i|} \nabla_{\theta} \log \pi_{\theta} (y_{i,t} | x, y_{i,<t}) 
\Bigg].
\label{equ:gspo_grad}
\end{align}

For comparison, the gradient of the GRPO objective is as follows (note that $\widehat{A}_{i,t} = \widehat{A}_{i}$):

\begin{align}
\nabla_{\theta} \mathcal{J}_\text{GRPO} (\theta) 
=&\ 
\nabla_{\theta} \mathbb{E}_{ x \sim \mathcal{D},\, \{y_i\}_{i=1}^G \sim \pi_{\theta_\text{old}}( \cdot | x) }
\left[ \frac{1}{G} \sum_{i=1}^{G} \frac{1}{|y_i|} \sum_{t=1}^{|y_i|} 
w_{i,t}(\theta) \widehat{A}_{i,t}
\right] \\
=&\
\mathbb{E}_{ x \sim \mathcal{D},\, \{y_i\}_{i=1}^G \sim \pi_{\theta_\text{old}}( \cdot | x) }
\left[ \frac{1}{G} \sum_{i=1}^{G} \widehat{A}_{i} 
\times \frac{1}{|y_i|} \sum_{t=1}^{|y_i|} 
\frac{ \pi_{\theta} (y_{i,t} | x, y_{i,<t}) }{ \pi_{\theta_\text{old}} (y_{i,t} | x,y_{i,<t})} 
\nabla_{\theta} \log \pi_{\theta} (y_{i,t} | x, y_{i,<t})  
\right].
\end{align}

Consequently, the main difference between GSPO and GRPO centers on \textit{the method by which each approach assigns weights to the gradients of token log likelihoods}. GRPO employs an ``importance weight'' for each token, represented as $\frac{ \pi_{\theta} (y_{i,t} | x, y_{i,<t}) }{ \pi_{\theta_\text{old}} (y_{i,t} | x,y_{i,<t})}$, so tokens contribute unequally according to this ratio. These varying weights—which can fall within the range $(0, 1+\varepsilon]$ when $\widehat{A}_{i} >0$ or within $[1-\varepsilon, +\infty)$ when $\widehat{A}_{i} < 0$—cannot be ignored and may have a compounding effect, potentially causing erratic behavior during optimization. By contrast, GSPO applies identical weights to every token within a response, thereby removing the source of instability introduced by the non-uniform weighting in GRPO.

\subsection{GSPO-token: A Token-level Objective Variant}

In contexts such as multi-turn RL, there may be a need for a more detailed adjustment of advantage beyond the sequence scale. To address this, we propose a token-level adaptation of GSPO, referred to as \textbf{GSPO-token}, which enables the customization of advantages on a per-token basis:

\begin{align}
\mathcal{J}_\text{GSPO-token}(\theta)
=
\mathbb{E}_{ x \sim \mathcal{D},\, \{y_i\}_{i=1}^G \sim \pi_{\theta_\text{old}}( \cdot | x) }
\left[ 
\frac{1}{G} \sum_{i=1}^{G} \frac{1}{|y_i|} \sum_{t=1}^{|y_i|} 
\min \left( s_{i,t}(\theta) \widehat{A}_{i,t},  \, \mathrm{clip} \left( s_{i,t}(\theta), 1 - {\varepsilon}, 1 + {\varepsilon} \right) \widehat{A}_{i,t} \right) 
\right],
\label{equ:gspo-token}
\end{align}

where

\begin{align}
s_{i,t}(\theta) 
= \mathrm{sg} \left[ s_{i}(\theta) \right]  \times \frac{ \pi_{\theta} (y_{i,t} | x, y_{i,<t}) }{ \mathrm{sg} \left[ \pi_{\theta} (y_{i,t} | x, y_{i,<t}) \right] },
\end{align}

and $\mathrm{sg}[\cdot]$ denotes only taking the numerical value but stopping the gradient, corresponding to the \texttt{detach} operation in PyTorch. The gradient of GSPO-token can be derived as:

\begin{align}
\nabla_{\theta} \mathcal{J}_\text{GSPO-token}(\theta)
=&\ 
\nabla_{\theta} \mathbb{E}_{x \sim \mathcal{D},\, \{y_i\}_{i=1}^G \sim \pi_{\theta_\text{old}}(\cdot | x)}
\left[ 
\frac{1}{G} \sum_{i=1}^{G}
\frac{1}{|y_i|} \sum_{t=1}^{|y_i|} 
s_{i,t}(\theta) \widehat{A}_{i,t}
\right] \\
=&\ 
\mathbb{E}_{x \sim \mathcal{D},\, \{y_i\}_{i=1}^G \sim \pi_{\theta_\text{old}}(\cdot | x)}
\left[ 
\frac{1}{G} \sum_{i=1}^{G} s_i(\theta) \cdot \frac{1}{|y_i|} \sum_{t=1}^{|y_i|} 
\widehat{A}_{i,t} \frac{\nabla_{\theta} \pi_{\theta}(y_{i,t} | x, y_{i,<t})}{\pi_{\theta}(y_{i,t} | x, y_{i,<t})}
\right] \\
=&\ 
\mathbb{E}_{x \sim \mathcal{D},\, \{y_i\}_{i=1}^G \sim \pi_{\theta_\text{old}}(\cdot | x)}
\left[ 
\frac{1}{G} \sum_{i=1}^{G} \left( \frac{\pi_{\theta}(y_i | x)}{\pi_{\theta_\text{old}}(y_i | x)} \right)^{\frac{1}{|y_i|}}
\cdot \frac{1}{|y_i|} \sum_{t=1}^{|y_i|}
\widehat{A}_{i,t} \nabla_{\theta} \log \pi_{\theta}(y_{i,t} | x, y_{i,<t})  
\right].
\label{equ:gspo-token_grad}
\end{align}

Observe that $\frac{ \pi_{\theta} (y_{i,t} | x, y_{i,<t}) }{ \mathrm{sg} \left[ \pi_{\theta} (y_{i,t} | x, y_{i,<t}) \right] }$ evaluates to 1 numerically; therefore, $s_{i,t}(\theta)$ is equal to $s_{i}(\theta)$ in terms of numeric value. Examining Equation~\eqref{equ:gspo} alongside~\eqref{equ:gspo-token}, and Equation~\eqref{equ:gspo_grad} with~\eqref{equ:gspo-token_grad}, we find that both GSPO-token and GSPO yield identical results for the objective, clipping criteria, and analytical gradient when assigning a constant advantage $\widehat{A}_{i,t} = \widehat{A}_{i}$ across all tokens within the response $y_i$. Nevertheless, the GSPO-token framework provides greater adaptability by enabling the assignment of distinct advantages to individual tokens.

\section{Experiments and Discussion}

\subsection{Empirical Results}

We conduct experiments using a cold-start model that is fine-tuned from Qwen3-30B-A3B-Base, evaluating both training reward trajectories and model performance metrics on several benchmarks: AIME'24 (reporting average Pass@1 across 32 samples), LiveCodeBench (202410-202502, using average Pass@1 from 8 samples), and CodeForces (utilizing the Elo Rating). Throughout the RL training process, each rollout batch is divided into four distinct mini-batches to facilitate gradient updates. For GSPO, we configure the left and right clipping bounds in Equation~\eqref{equ:gspo} to be 3e-4 and 4e-4, respectively. Our baseline for comparison is GRPO, for which we set the left and right clipping thresholds in Equation~\eqref{equ:grpo} to 0.2 and 0.27; these values are selected following extensive hyperparameter tuning to guarantee a fair assessment. It is important to mention that GRPO requires the Routing Replay training scheme to achieve reliable convergence for MoE RL, a topic further detailed in \S~\ref{subsec:routing_replay}. In contrast, \textit{GSPO eliminates the necessity of this training strategy}.

Figure~\ref{fig:results} illustrates that GSPO ensures stable progress during training. Notably, \textbf{GSPO consistently enhances performance as training compute increases, the query set is periodically refreshed, and the generation length is extended}. In addition, GSPO outperforms GRPO in terms of training efficiency, achieving higher training accuracy and superior benchmark results given equal training compute and query usage. Lastly, we have effectively utilized GSPO for RL training on recent Qwen3 models, which provides compelling evidence of GSPO’s effectiveness in leveraging RL scaling benefits for large language models.

\subsection{Curious Observation on Clipping Fractions}

An important difference between GSPO and GRPO is that GSPO performs clipping at the sequence level, affecting whole responses, while GRPO clips at the token level. As illustrated in Figure~\ref{fig:clipping}, this leads to a gap of approximately two orders of magnitude in the proportion of tokens clipped by GSPO compared to GRPO; modifying the clipping ranges does not affect this pronounced difference. Interestingly, although GSPO ends up clipping many more tokens—thereby relying on a smaller subset of tokens for training or gradient estimation—it nonetheless attains greater training efficiency than GRPO. This somewhat surprising outcome, where extensive token clipping correlates with improved training efficiency, suggests that GRPO's approach of using token-level gradient estimates introduces substantial noise and diminishes the effectiveness of sample exploitation. By contrast, the sequence-level methodology in GSPO yields a more robust and informative learning signal.

\subsection{Benefit of GSPO for MoE Training}
\label{subsec:routing_replay}

\paragraph{Background}
In contrast to RL training for dense models, the sparsely activated architecture of MoE models gives rise to distinct stability problems. Notably, when employing the GRPO algorithm, the \textit{expert-activation volatility} observed in MoE models can impede effective convergence during RL training. Specifically, we observed that the set of experts triggered for an identical response can vary substantially following one or several gradient update steps. To illustrate, using the 48-layer Qwen3-30B-A3B-Base configuration, each RL gradient update causes about 10\% of the experts activated by the updated policy $\pi_{\theta}$, for a given rollout sample, to differ from those chosen by the preceding policy $\pi_{\theta_\text{old}}$.
This issue intensifies in deeper MoE architectures, leading to significant variability in the token-level importance ratios $w_{i,t}(\theta) = \frac{ \pi_{\theta} (y_{i,t} | x, y_{i,<t}) }{ \pi_{\theta_\text{old}} (y_{i,t} | x,y_{i,<t})}$, rendering these ratios unreliable, as detailed in \S~\ref{sec:motivation} and~\ref{subsec:gradient}. As a result, this undermines the normal convergence behavior of RL optimization.

\paragraph{Our Previous Approach} In our prior work, we addressed this issue by utilizing the \textbf{Routing Replay} method during training. In detail, we store the experts chosen by $\pi_{\theta_\text{old}}$ and later use these stored routing configurations within $\pi_{\theta}$ during the evaluation of importance ratios, specifically $w_{i,t}(\theta) = \frac{ \pi_{\theta} (y_{i,t} | x, y_{i,<t}) }{ \pi_{\theta_\text{old}} (y_{i,t} | x,y_{i,<t})}$. This ensures that, for each token $y_{i,t}$, the computation of both $\pi_{\theta} (y_{i,t} | x, y_{i,<t})$ and $\pi_{\theta_\text{old}} (y_{i,t} | x,y_{i,<t})$ is carried out with the same set of experts activated. As a result, this technique mitigates volatility in the calculated importance ratios at the token level and allows optimization to proceed smoothly over a stable, consistently activated subnetwork throughout gradient steps. Figure~\ref{fig:routing_replay} illustrates the critical role Routing Replay plays in promoting stable convergence during the GRPO-based training of MoE architectures.

\paragraph{Benefit of GSPO} While Routing Replay allows GRPO training for MoE models to achieve convergence, this approach entails extra memory and communication requirements due to the repeated use of routing patterns, which can in turn restrict the model’s effective capacity. On the other hand, as illustrated in Figure~\ref{fig:results}, GSPO removes the reliance on Routing Replay. With GSPO, the computation of the importance ratios $s_i(\theta)$ can proceed in a standard fashion, achieving both stable convergence and optimization. The principal advantage of GSPO is its exclusive focus on sequence-level likelihoods (i.e., $\pi_{\theta} (y_i | x)$), rather than the likelihood of individual tokens (i.e., $\pi_{\theta} (y_{i,t} | x, y_{i,<t})$), leading to greater robustness. Because the language modeling competence of the MoE architecture is preserved, the overall sequence likelihood remains comparatively stable throughout training. Overall, GSPO directly addresses the problem of instability in expert activation for MoE models, removing the necessity for elaborate solutions such as Routing Replay. This streamlines the training workflow, enhances stability, and permits exploitation of the model’s full representational capacity.

\subsection{Benefit of GSPO for RL Infrastructure}

Due to the variation in numerical precision between training frameworks (such as Megatron) and inference frameworks (like SGLang and vLLM), practitioners generally rely on the training engine to recalculate the likelihoods of generated outputs according to the previous policy $\pi_{\theta_\text{old}}$. In contrast, GSPO focuses on optimizing sequence-level rather than token-level likelihoods, which are inherently less sensitive to discrepancies in computational precision. Consequently, GSPO enables the direct utilization of likelihoods provided by the inference engine for optimization purposes, eliminating the need to recompute them with the training engine. This feature is particularly advantageous in contexts such as partial rollout, multi-turn RL, and systems that separate the training and inference processes.

\section{Conclusion}

We introduce Group Sequence Policy Optimization (GSPO), a novel reinforcement learning algorithm specifically designed for the training of large language models. Grounded in the fundamental idea of importance sampling, GSPO establishes importance ratios using sequence likelihoods and incorporates sequence-level techniques such as clipping, reward assignment, and optimization. Compared to GRPO, GSPO delivers significantly enhanced stability during training, higher efficiency, and superior performance. Its effectiveness is especially pronounced in large-scale reinforcement learning applications with MoE models, underpinning the remarkable advancements observed in the latest Qwen3 models. As a scalable foundation, GSPO enables continued expansion of reinforcement learning approaches, paving the way for substantial progress in artificial intelligence.

\end{document}